\documentclass[a4paper]{article}
\usepackage{amsfonts}
\usepackage{amsmath}
\usepackage{amssymb}
\usepackage{graphicx}     

\begin{document}
  
\noindent \large Auteur: Nathan De Roy \\
\noindent \large Studierichting: Informatica\\
\noindent \large Oefeningenreeks 1C nr. 12.1\\

\medskip

\normalsize

Los $\dfrac{-1 + i}{\sqrt{2}}$ op\\

$z_1$ = $-1 + i$ = $r_1 \operatorname{cis}(\alpha)$\\

en $z_2$ = $r_2 \operatorname{cis}(\beta)$\\

Moduli:\\

$r_1$ = $\sqrt{(-1)^2 + 1^2}$ = $\sqrt{2}$\\

$r_2$ = $\sqrt{2}$\\

Argumenten:\\

$\alpha$ = $\tan^{-1}\left(\dfrac{b}{a}\right)$ = $\tan^{-1}\left(\dfrac{1}{-1}\right)$ = $\tan^{-1}(-1)$ = $\dfrac{3 \pi}{4}$\\

$\beta$ = $\tan^{-1}\left(\dfrac{0}{\sqrt{2}}\right)$\\

Goniometrische vorm:\\

$z_1$ = $\sqrt{2} \operatorname{cis}(\dfrac{3 \pi}{4})$\\

$z_2$ = $\sqrt{2} \operatorname{cis}(0)$\\

Deling:\\

$\dfrac{z_1}{z_2}$ = $\dfrac{z_1}{r_2}\operatorname{cis}(\alpha-\beta)$ = $\operatorname{cis}\left(\dfrac{3\pi}{4}\right)$\\

\begin{thebibliography}{999}
%\bibitem{a} Referentie
\end{thebibliography}
\end{document} 
